\documentclass[11pt, a4paper]{article}
\usepackage[a4paper, margin=2.5cm]{geometry}
\usepackage[utf8]{inputenc}
\usepackage[T1]{fontenc}
\usepackage[spanish]{babel}
\usepackage{amsmath, amssymb, booktabs}
\usepackage{graphicx}
\usepackage{hyperref}
\usepackage{xcolor}
\usepackage{listings}

\hypersetup{
    colorlinks=true,
    linkcolor=blue!50!black,
    urlcolor=blue!50!black,
    citecolor=blue!50!black
}

\lstset{
    basicstyle=\ttfamily\small,
    breaklines=true,
    frame=single,
    backgroundcolor=\color{gray!5},
    rulecolor=\color{gray!30}
}

\title{\textbf{Informe Técnico: Simulador de Datos Sísmicos (SEG-Y) y Optimización de E/S Mediante Estructuras de Datos y Prefetching}}
\author{Proyecto de Arquitectura de Computadores / Estructuras de Datos}
\author{
    \begin{tabular}{c}
        \textbf{Jose Alvarez} \\ C.I.: 31389258
    \end{tabular}
    \hspace{2cm}
    \begin{tabular}{c}
        \textbf{Jose Nuñez} \\ C.I.: 31581050
    \end{tabular}
    \\[1.5em]
    \textit{Universidad de Carabobo} \\
    \textit{Facultad de Ciencias y Tecnología} \\
    \textit{Departamento de Computación -- Arquitectura del Computador}
}
\date{Octubre 2026}

\begin{document}

\maketitle

\begin{abstract}
En este informe se presenta el desarrollo, implementación y análisis de rendimiento de un simulador de datos sísmicos en formato SEG-Y desarrollado en C++17. Se abordan los desafíos asociados al manejo de volúmenes de datos masivos en geofísica, evaluando diferentes backends de almacenamiento en memoria, indexación espacial basada en árboles rojinegros y estrategias de carga anticipada (\textit{prefetching}) con caché LRU para mitigar los cuellos de botella de Entrada/Salida (E/S). Los resultados demuestran que el uso de prefetching espacial incrementa la tasa de aciertos (\textit{hit ratio}) al 100\% y reduce drásticamente las operaciones de lectura física.
\end{abstract}

\section{El Problema}
En la industria geofísica y de exploración petrolera, los conjuntos de datos sísmicos tridimensionales (3D) representan volúmenes masivos de información estructurada que frecuentemente alcanzan gigabytes o terabytes. Un archivo en formato estándar \textit{SEG-Y} contiene cabeceras globales y miles de trazas individuales (registros de ondas reflejadas en el subsuelo). 

El problema fundamental radica en que los patrones de acceso durante la interpretación interactiva o el procesamiento sísmico suelen ser complejos (secuenciales, inversos, aleatorios o con saltos espaciales). Cuando un software accede a estos volúmenes sin una estrategia adecuada de gestión de memoria y caché, se producen severos cuellos de botella de Entrada/Salida (E/S) y una saturación ineficiente de la memoria RAM. El objetivo principal de este proyecto es construir un banco de pruebas que permita simular, medir y comparar el impacto de distintas arquitecturas de software y algoritmos de optimización sobre un volumen sísmico sintético.

\section{Técnicas Computacionales para su Solución}
Para resolver los desafíos de rendimiento y almacenamiento, el simulador se estructuró de manera modular en C++17, implementando los siguientes componentes algorítmicos y estructuras de datos:

\subsection{Parser del Formato SEG-Y (\texttt{SEGYParser})}
El sistema implementa un analizador robusto capaz de leer y procesar la cabecera binaria de 400 bytes (para extraer parámetros como el intervalo de muestreo temporal $\Delta t$ y el formato de codificación de las muestras) y recorrer iterativamente cada traza con su respectiva cabecera de 240 bytes. Asimismo, incluye un generador sintético basado en ondas \textit{Ricker} correlacionadas espacialmente que permite crear cubos de prueba configurables (por ejemplo, dimensiones de $64 \times 64$ o $128 \times 128$ trazas con 256/512 muestras por traza).

\subsection{Backends de Almacenamiento en Memoria (\texttt{StorageBackend})}
Se diseñó una interfaz abstracta para comparar dos estructuras de contenedores de la STL:
\begin{itemize}
    \item \textbf{\texttt{VectorStorage} ($\text{std::vector}$):} Almacena las trazas en un bloque de memoria contiguo, maximizando la localidad espacial y optimizando los accesos aleatorios, a costo de posibles reasignaciones dinámicas si el tamaño varía.
    \item \textbf{\texttt{DequeStorage} ($\text{std::deque}$):} Utiliza bloques de memoria fragmentados pero encadenados, ofreciendo un rendimiento eficiente en inserciones masivas durante la carga inicial del archivo.
\end{itemize}

\subsection{Indexación Espacial (\texttt{TraceIndex})}
En geofísica, las trazas no se consultan por su índice físico secuencial, sino por sus coordenadas lógicas en la malla de adquisición: el par de coordenadas \textit{(inline, crossline)}. Para resolver consultas espaciales en tiempo óptimo, se implementó un índice basado en un árbol rojo-negro a través de $\text{std::map}$. Este índice relaciona la clave lógica \textit{(inline, crossline)} con la posición física exacta del vector de almacenamiento en un tiempo logarítmico $\mathcal{O}(\log n)$.

\subsection{Gestión de Caché y Prefetching (\texttt{Prefetcher})}
Para evitar accesos repetitivos al almacenamiento, se desarrolló un sistema de caché con política de reemplazo \textit{LRU} (\textit{Least Recently Used}) acoplado a un algoritmo de \textit{prefetching} espacial. Cuando una traza $i$ es solicitada, el prefetcher predice y carga de forma proactiva un número configurable de trazas vecinas espaciales ($i+1, i+2, \dots, i+k$), anticipándose al comportamiento de lectura del usuario y reduciendo la latencia de E/S.

\section{Análisis y Comparación de Resultados}
El simulador fue ejecutado sobre un archivo SEG-Y sintético de prueba ($\sim 4.94\text{ MB}$, compuesto por 4096 trazas de 256 muestras cada una). A continuación se analizan los resultados empíricos obtenidos en el banco de pruebas:

\subsection{Rendimiento de Carga y Uso de Memoria}
La carga del volumen completo demostró ser sumamente ligera en ambas estructuras de la STL, operando con un consumo de memoria RAM aproximado de $4.25\text{ MB}$:
\begin{itemize}
    \item \textbf{$\text{std::vector}$:}{ Tiempo de carga de $57.0\text{ ms}$.}
    \item \textbf{$\text{std::deque}$:}{ Tiempo de carga de $70.0\text{ ms}$.}
\end{itemize}
Por su parte, la construcción del árbol rojo-negro para indexar las 4096 trazas tomó apenas $1.0\text{ ms}$, consumiendo una huella de memoria auxiliar de $192.00\text{ KB}$.

\subsection{Impacto del Prefetching en Patrones de Acceso}
La evaluación bajo diferentes patrones de consulta (secuencial, inverso, aleatorio y con saltos) evidenció el comportamiento crítico de la caché. Sin prefetching, el \textit{Hit Ratio} se mantiene en $0.000$, obligando al sistema a realizar la totalidad de las 4096 cargas físicas al backend en cada recorrido.

Sin embargo, al activar el mecanismo de \textit{prefetching} con profundidades de 2, 4 u 8 vecinos bajo un patrón secuencial, se observa el impacto analítico clave del proyecto:

\begin{table}[h]
    \centering
    \caption{Rendimiento del sistema bajo diferentes niveles de Prefetching (Patrón Secuencial).}
    \label{tab:prefetch_results}
    \begin{tabular}{lrrrr}
        \toprule
        \textbf{Prefetch (Vecinos)} & \textbf{Tiempo (ms)} & \textbf{Hit Ratio} & \textbf{Cargas Físicas} & \textbf{Prefetches} \\
        \midrule
        off (Sin caché) & $2.0$ & $0.000$ & $4096$ & -- \\
        2 vecinos       & $10.0$ & $1.000$ & $1$ & -- \\
        4 vecinos       & $5.0$ & $1.000$ & $1$ & -- \\
        8 vecinos       & $4.0$ & $1.000$ & $1$ & -- \\
        \bottomrule
    \end{tabular}
\end{table}

Como se aprecia en la Tabla \ref{tab:prefetch_results}, al habilitar la precarga espacial, la tasa de aciertos (\textit{Hit Ratio}) salta inmediatamente a \textbf{$1.000$ (100\% de aciertos)} y el número de accesos físicos al almacenamiento se desploma drásticamente de 4096 operaciones a \textbf{tan solo 1 carga inicial}. Esto demuestra empíricamente que la predicción basada en la localidad espacial mitiga de forma radical los cuellos de botella en el procesamiento de datos sísmicos.

\section{Conclusiones}
El desarrollo e implementación de este simulador permitió comprobar que la elección adecuada de estructuras de datos contiguas ($\text{std::vector}$), combinada con una indexación espacial eficiente mediante árboles rojinegros y políticas de caché con \textit{prefetching}, optimiza de manera exponencial el rendimiento en la lectura de cubos sísmicos SEG-Y. El proyecto cumple a cabalidad con los requerimientos de arquitectura de alto rendimiento y provee una base sólida para futuras extensiones en entornos distribuidos o de almacenamiento en disco secundario.

\end{document}